% Generated draft: requires mathematical and source review.
\documentclass[11pt]{article}
\usepackage[margin=25mm]{geometry}
\usepackage{fontspec}
\setmainfont{Latin Modern Roman}
\usepackage{amsmath,amssymb,mathtools,unicode-math}
\setmathfont{Latin Modern Math}
\usepackage{enumitem,needspace,xurl,hyperref}
\setlength{\parindent}{0pt}
\setlength{\parskip}{5pt}
\setlength{\emergencystretch}{3em}
\newcommand{\N}{\mathbb{N}}
\newcommand{\Z}{\mathbb{Z}}
\newcommand{\Q}{\mathbb{Q}}
\newcommand{\R}{\mathbb{R}}
\newcommand{\C}{\mathbb{C}}
\newcommand{\F}{\mathbb{F}}
\newcommand{\A}{\mathbb{A}}
\newcommand{\PP}{\mathbb{P}}
\newcommand{\E}{\mathbb{E}}
\newcommand{\eps}{\varepsilon}
\newcommand{\dd}{\,\mathrm{d}}
\begin{document}
\section*{UnsolvedMath ID: 10800007}
\textbf{Source problem code:} catalogue entry
\subsection{Definitions and mathematical statement}
Let \(p=(a,b)\in\R^6\) denote the source's parameter pair and let \(f_a,g_b\) be its two real functions on \(\R^2\). Their formulas, degrees, and parameter split are not included in the supplied record. Define the solution incidence space, conditional on that data, by
\[
 E=\{(p,x)\in\R^6\times\R^2:f_a(x)=0,\ g_b(x)=0\},
 \qquad\pi:E\to\R^6,\quad\pi(p,x)=p.
\]
A continuous choice of a solution on an open set \(U\subseteq\R^6\) is a map \(\varphi:U\to\R^2\) with
\[
 f_a(\varphi(a,b))=g_b(\varphi(a,b))=0
 \quad\text{for every }(a,b)\in U.
\]
Equivalently, \(p\mapsto(p,\varphi(p))\) is a continuous section of \(\pi\) over \(U\).

\emph{Problem (Victor Vassiliev, Problem 2A).} Determine the minimum number of open sets covering \(\R^6\) on each of which such a continuous solution choice exists. If the minimum is finite, it is the least positive integer \(m\) such that there are open \(U_1,\ldots,U_m\) and continuous \(\varphi_i:U_i\to\R^2\) satisfying
\[
 \bigcup_{i=1}^{m}U_i=\R^6,\qquad
 f_a(\varphi_i(a,b))=g_b(\varphi_i(a,b))=0
 \quad((a,b)\in U_i).
\]
The source does not remove singular parameters or require the chosen solutions to be unique, distinct, or regular. The maps need not agree on overlaps. The formulas of \(f_a,g_b\) are indispensable: without them, even solvability of every fiber and existence of local sections cannot be checked. The excerpt also does not establish finiteness of the requested covering number.
\subsection{Short English statement}
Find the smallest open cover of the full parameter space \(\R^6\) on whose members one can continuously choose a common real zero of \(f_a\) and \(g_b\). The defining functions are missing, so the associated solution-selection problem is not yet fully specified.
\subsection{Sources}
\begin{itemize}
\item Local UnsolvedMath registry, permanent ID 10800007, source code catalogue entry; proposer Victor Vassiliev, Problem 2A.
\item Vassiliev, \emph{A Few Problems on Monodromy and Discriminants}, Problem 2A. Supplied URL: \url{https://amj.math.stonybrook.edu/html-articles/Files-2015-2024/15-11/}; recorded access 2026-07-29; extraction html-structure. The article contains eighteen explicitly labeled theorem-style problem environments.
\item Supplied status open and OPEN-TRIAGE. The research summary says a lower bound of \(2\) is known and it did not determine whether this map's covering number, or higher-degree analogues, equals \(2\) or is larger. The specific map definition and supporting lower-bound reference are not supplied. Current primary-literature status and redistribution terms of the publicly accessible author/publisher copy remain NEEDS\_REVIEW.
\end{itemize}
\end{document}
